% GENERATED FILE. Edit canonical lessons, then run npm run generate:books.
% canonical-source-hash: fnv1a64:b71710ead85f7004
% canonical-lessons: GE-C283-zeuge, GE-C283-aussagen, GE-C283-aussage, GE-C283-beobachten, GE-C283-schildern

\chapter{Der Zeuge}
\label{ch:ge-der-zeuge}
\hlchaptermodality{283}

\begin{chapteropening}
\textbf{By the end of this chapter:} \emph{I can talk about witnesses: a witness, to testify, a statement, to observe and to recount.}

Complete the last lesson of chapter 283: I can talk about witnesses: a witness, to testify, a statement, to observe and to recount.
\end{chapteropening}

\section[der Zeuge]{der Zeuge — the witness}

\label{lesson:GE-C283-zeuge}



\begin{quote}
Before the new one: say the German for to threaten, then the German for the weapon.
\end{quote}

\subsection*{You'll want to know: der Zeuge}
\textbf{der Zeuge} — \textquotedblleft{}the witness\textquotedblright{}.

\begin{quote}
Ein Zeuge hatte alles gesehen.
\end{quote}

\begin{quote}
\emph{A witness had seen everything.}
\end{quote}

Weak noun: \emph{den Zeugen}, \emph{dem Zeugen}. Feminine \emph{die Zeugin}.

\begin{cousinweb}
\emph{Zeuge} is usually traced to Middle High German \emph{geziuc} \textquotedblleft{}evidence\textquotedblright{}, built on the root of \emph{ziehen} \textquotedblleft{}to pull\textquotedblright{}: someone called in. EtymDB lists only its Dutch cousin \emph{getuige}.
\end{cousinweb}

\begin{grammarlens}[title={What you've built}]
One more word for telling, explaining and describing at B1.
\end{grammarlens}

\subsection*{Your turn}
\begin{itemize}
\raggedright
  \item \emph{Say it:} \emph{der Zeuge}
  \item \emph{Say it:} \emph{der Zeuge}, once more
  \item \emph{Say it:} a sentence of your own with \emph{der Zeuge}
  \item \emph{Recall:} say the German for to threaten, then the German for the weapon, then say \emph{der Zeuge} again
\end{itemize}

\begin{tcolorbox}[breakable,colback=teal!4,colframe=teal!35!black,title={Before you move on}]
What does der Zeuge mean? (\textquotedblleft{}The witness\textquotedblright{}.) Say it once more.
\end{tcolorbox}
\section[aussagen]{aussagen — to testify, to give evidence}

\label{lesson:GE-C283-aussagen}



\begin{quote}
Before the new one: say the German for the weapon, then the German for the witness.
\end{quote}

\subsection*{You'll want to know: aussagen}
\textbf{aussagen} — \textquotedblleft{}to testify, to give evidence\textquotedblright{}.

\begin{quote}
Die Nachbarin sagte vor Gericht aus.
\end{quote}

\begin{quote}
\emph{The neighbour testified in court.}
\end{quote}

Separable: \emph{sie sagt aus}, \emph{sie hat ausgesagt}.

\begin{cousinweb}
\emph{aussagen} joins \emph{aus} \textquotedblleft{}out\textquotedblright{} and \emph{sagen} \textquotedblleft{}to say\textquotedblright{}, an inherited Germanic verb with the Dutch cousin \emph{zeggen} (EtymDB): to say out.
\end{cousinweb}

\begin{grammarlens}[title={What you've built}]
One more word for telling, explaining and describing at B1.
\end{grammarlens}

\subsection*{Your turn}
\begin{itemize}
\raggedright
  \item \emph{Say it:} \emph{aussagen}
  \item \emph{Say it:} \emph{aussagen}, once more
  \item \emph{Say it:} a sentence of your own with \emph{aussagen}
  \item \emph{Recall:} say the German for the weapon, then the German for the witness, then say \emph{aussagen} again
\end{itemize}

\begin{tcolorbox}[breakable,colback=teal!4,colframe=teal!35!black,title={Before you move on}]
What does aussagen mean? (\textquotedblleft{}To testify, to give evidence\textquotedblright{}.) Say it once more.
\end{tcolorbox}
\section[die Aussage]{die Aussage — the statement, the testimony}

\label{lesson:GE-C283-aussage}



\begin{quote}
Before the new one: say the German for the witness, then the German for to testify.
\end{quote}

\subsection*{You'll want to know: die Aussage}
\textbf{die Aussage} — \textquotedblleft{}the statement, the testimony\textquotedblright{}.

\begin{quote}
Seine Aussage war sehr genau.
\end{quote}

\begin{quote}
\emph{His statement was very precise.}
\end{quote}

Plural: \emph{die Aussagen}.

\begin{cousinweb}
\emph{Aussage} is built on \emph{aussagen} \textquotedblleft{}to testify\textquotedblright{}, from \emph{sagen} \textquotedblleft{}to say\textquotedblright{}, an inherited Germanic verb (EtymDB).
\end{cousinweb}

\begin{grammarlens}[title={What you've built}]
One more word for telling, explaining and describing at B1.
\end{grammarlens}

\subsection*{Your turn}
\begin{itemize}
\raggedright
  \item \emph{Say it:} \emph{die Aussage}
  \item \emph{Say it:} \emph{die Aussage}, once more
  \item \emph{Say it:} a sentence of your own with \emph{die Aussage}
  \item \emph{Recall:} say the German for the witness, then the German for to testify, then say \emph{die Aussage} again
\end{itemize}

\begin{tcolorbox}[breakable,colback=teal!4,colframe=teal!35!black,title={Before you move on}]
What does die Aussage mean? (\textquotedblleft{}The statement, the testimony\textquotedblright{}.) Say it once more.
\end{tcolorbox}
\section[beobachten]{beobachten — to observe, to watch}

\label{lesson:GE-C283-beobachten}



\begin{quote}
Before the new one: say the German for to testify, then the German for the statement.
\end{quote}

\subsection*{You'll want to know: beobachten}
\textbf{beobachten} — \textquotedblleft{}to observe, to watch\textquotedblright{}.

\begin{quote}
Sie beobachtete den Mann vom Fenster aus.
\end{quote}

\begin{quote}
\emph{She watched the man from the window.}
\end{quote}

Regular: \emph{beobachtete}, \emph{hat beobachtet}.

\begin{cousinweb}
\emph{beobachten} is usually traced to a Low German verb built on \emph{achten} \textquotedblleft{}to pay heed\textquotedblright{}, from Old High German \emph{ahtōn} (EtymDB): to keep an eye on.
\end{cousinweb}

\begin{grammarlens}[title={What you've built}]
One more word for telling, explaining and describing at B1.
\end{grammarlens}

\subsection*{Your turn}
\begin{itemize}
\raggedright
  \item \emph{Say it:} \emph{beobachten}
  \item \emph{Say it:} \emph{beobachten}, once more
  \item \emph{Say it:} a sentence of your own with \emph{beobachten}
  \item \emph{Recall:} say the German for to testify, then the German for the statement, then say \emph{beobachten} again
\end{itemize}

\begin{tcolorbox}[breakable,colback=teal!4,colframe=teal!35!black,title={Before you move on}]
What does beobachten mean? (\textquotedblleft{}To observe, to watch\textquotedblright{}.) Say it once more.
\end{tcolorbox}
\section[schildern]{schildern — to recount, to describe vividly}

\label{lesson:GE-C283-schildern}



\begin{quote}
Before the new one: say the German for the statement, then the German for to observe.
\end{quote}

\subsection*{You'll want to know: schildern}
\textbf{schildern} — \textquotedblleft{}to recount, to describe vividly\textquotedblright{}.

\begin{quote}
Der Zeuge schilderte den Unfall genau.
\end{quote}

\begin{quote}
\emph{The witness recounted the accident precisely.}
\end{quote}

Stronger than \emph{beschreiben}: you paint the scene in words.

\begin{cousinweb}
\emph{schildern} is built on \emph{Schild} \textquotedblleft{}shield\textquotedblright{}, a cousin of English \emph{shield} (EtymDB): it first meant painting coats of arms, then painting a scene in words.
\end{cousinweb}

\begin{grammarlens}[title={What you've built}]
One more word for telling, explaining and describing at B1.
\end{grammarlens}

\subsection*{Your turn}
\begin{itemize}
\raggedright
  \item \emph{Say it:} \emph{schildern}
  \item \emph{Say it:} \emph{schildern}, once more
  \item \emph{Say it:} a sentence of your own with \emph{schildern}
  \item \emph{Recall:} say the German for the statement, then the German for to observe, then say \emph{schildern} again
\end{itemize}

\begin{tcolorbox}[breakable,colback=teal!4,colframe=teal!35!black,title={Before you move on}]
What does schildern mean? (\textquotedblleft{}To recount, to describe vividly\textquotedblright{}.) Say it once more.
\end{tcolorbox}
